\chapter{Operators and Arithmetic}
\label{ch:operators}

\begin{goals}
\begin{itemize}
  \item The arithmetic operators and the order they are applied in
  \item The comparison operators and their boolean results
  \item Combining conditions with \code{and}, \code{or} and \code{not}
  \item The short forms of assignment
  \item A few everyday calculations with these tools
\end{itemize}
\end{goals}

\section{What is an operator?}

The signs such as \code{+} and \code{*} that we calculate with are called
\textbf{operators}, and the things they work on are called
\textbf{operands}. In \code{2 + 3}, the sign \code{+} is the operator and
\code{2} and \code{3} are the operands. You have met a few operators
already; in this chapter we look at all the basic operators of Salam in one
place.

\section{Arithmetic operators}

\begin{center}
\begin{tabular}{clc}
\toprule
Operator & Meaning & Example \\
\midrule
\code{+} & addition & \code{7 + 2} is \code{9} \\
\code{-} & subtraction & \code{7 - 2} is \code{5} \\
\code{*} & multiplication & \code{7 * 2} is \code{14} \\
\code{/} & division & \code{7 / 2} is \code{3} \\
\code{\%} & remainder & \code{7 \% 2} is \code{1} \\
\code{**} & power & \code{2 ** 3} is \code{8} \\
\bottomrule
\end{tabular}
\end{center}

\program{The arithmetic operators}
\begin{salamcode}
func main:
    println 7 + 2, 7 - 2, 7 * 2
    println 7 / 2, 7 % 2
    println 2 ** 3, 2 ** 10
    println 10 - 15
end
\end{salamcode}

\begin{salamout}
9 5 14
3 1
8 1024
-5
\end{salamout}

The power operator, \code{**}, always gives a floating-point result, even
when both operands are integers. A minus sign can be placed before a number
or a variable to give its negative: \code{-age}.

\section{The order of operations}

When several operators appear in one expression, Salam follows the same rule
you learned in school mathematics: powers first, then multiplication,
division and remainder, then addition and subtraction. Operators of the same
rank are applied in the order they appear in the expression, one after the
other. The only exception is the power operator, which is worked out from
right to left: \code{2 ** 2 ** 3} means \code{2 ** (2 ** 3)}, and its value
is 256, not 64. The example below shows the order of the other operators:

\program{The order of operators}
\begin{salamcode}
func main:
    println 2 + 3 * 4
    println (2 + 3) * 4
    println 10 - 4 - 3
    println 100 / 10 / 2
    println 2 * 3 ** 2
end
\end{salamcode}

\begin{salamout}
14
20
3
5
18
\end{salamout}

\code{2 + 3 * 4} came to \code{14}, not \code{20}, because multiplication is
done first. \code{10 - 4 - 3} is \code{3}, because \code{10 - 4} is worked
out first and then \code{6 - 3}. Wherever you are in doubt, or want a
different order, use parentheses. An extra pair of parentheses never does
any harm, and it often makes the expression easier to read.

\begin{tip}
Break very long and complicated expressions into several variables with
descriptive names. Instead of writing \code{println (a + b) * c / 100 - d}
in one go, first write \code{sum := a + b}, then work out
\code{discount := sum * c / 100}, and finally use \code{println discount - d}.
The reader (and you yourself a week later) will thank you.
\end{tip}

\section{What division and remainder can do}

Integer division and remainder, together, are a powerful tool for ``taking
numbers apart''. A few examples:

\program{Converting seconds to hours and minutes}
\begin{salamcode}
func main:
    total seconds := 3725
    hours := total seconds / 3600
    minutes := (total seconds % 3600) / 60
    seconds := total seconds % 60
    println hours, "hours,", minutes, "minutes and", seconds, "seconds"
end
\end{salamcode}

\begin{salamout}
1 hours, 2 minutes and 5 seconds
\end{salamout}

Every hour is 3600 seconds, so integer division by 3600 gives the number of
whole hours. The remainder of that division is the seconds that did not make
up a whole hour, and we divide those by 60 to get the minutes.

\program{The digits of a number}
\begin{salamcode}
func main:
    number := 472
    println "Ones:", number % 10
    println "Tens:", (number / 10) % 10
    println "Hundreds:", number / 100
end
\end{salamcode}

\begin{salamout}
Ones: 2
Tens: 7
Hundreds: 4
\end{salamout}

The remainder of division by 10 is always the ones digit. Dividing by 10
drops one digit from the right-hand end. By combining the two you can reach
any digit.

\program{Percentages and discounts}
\begin{salamcode}
func main:
    price := 240
    discount percent := 15
    discount := (price * discount percent) / 100
    println "Discount:", discount
    println "Final price:", price - discount
end
\end{salamcode}

\begin{salamout}
Discount: 36
Final price: 204
\end{salamout}

Notice that we multiplied first and divided afterwards. Had we written
\code{price * (discount percent / 100)}, the integer division of 15 by 100
would have been zero and the whole discount would have vanished. With
integers, do the division as late as you can.

\section{Short forms of assignment}

In chapter \ref{ch:variables} you met \code{+=} and \code{++}. The other
arithmetic operators have short forms too:

\begin{center}
\begin{tabular}{ll}
\toprule
Short form & Equivalent \\
\midrule
\code{a += 5} & \code{a = a + 5} \\
\code{a -= 5} & \code{a = a - 5} \\
\code{a *= 5} & \code{a = a * 5} \\
\code{a /= 5} & \code{a = a / 5} \\
\code{a \%= 5} & \code{a = a \% 5} \\
\code{a++} & \code{a = a + 1} \\
\code{a--} & \code{a = a - 1} \\
\bottomrule
\end{tabular}
\end{center}

The power operator has a short form \code{**=} as well, but since the result
of a power is always floating-point, it can only be used on a variable of
type \code{f64}.

\program{The short forms}
\begin{salamcode}
func main:
    mut score := 100
    score += 20
    println score
    score -= 50
    println score
    score *= 2
    println score
    score /= 4
    println score
    score++
    println score
    score--
    println score
end
\end{salamcode}

\begin{salamout}
120
70
140
35
36
35
\end{salamout}

\section{Comparison operators}

The comparison operators weigh two values against each other and give a
boolean result: \code{true} or \code{false}.

\begin{center}
\begin{tabular}{ccl}
\toprule
Operator & Word & Meaning \\
\midrule
\code{==} & \code{eq} & is equal to \\
\code{!=} & \code{neq} & is not equal to \\
\code{<} & \code{lt} & is less than \\
\code{>} & \code{gt} & is greater than \\
\code{<=} & \code{lte} & is less than or equal to \\
\code{>=} & \code{gte} & is greater than or equal to \\
\bottomrule
\end{tabular}
\end{center}

\program{Comparisons}
\begin{salamcode}
func main:
    age := 20
    println age == 20
    println age != 20
    println age < 18
    println age >= 18
    println "apple" == "apple"
    println "apple" == "pear"
end
\end{salamcode}

\begin{salamout}
true
false
false
true
true
false
\end{salamout}

\begin{tip}
In place of each comparison operator you may write the word beside it in the
table above: \code{println age eq 20} does exactly what
\code{println age == 20} does, and \code{println age gte 18} is the same as
\code{println age >= 18}. In this book we mostly use the signs.
\end{tip}

\begin{warn}
Confusing \code{=} with \code{==} is one of the most common mistakes made by
programmers everywhere. \code{=} assigns a value; \code{==} asks ``are
they equal?''. If you write \code{=} where a comparison belongs, Salam
reports an error, which is good news.
\end{warn}

The result of a comparison is an ordinary boolean value and can be kept in a
variable. When a condition is complicated or is used several times, this
makes the program easier to read:

\begin{salamcode}
func main:
    temperature := 38
    is hot := temperature > 30
    println "Is it hot?", is hot
end
\end{salamcode}

\begin{salamout}
Is it hot? true
\end{salamout}

\section{Combining conditions: and, or, not}

Many decisions depend on more than one condition: ``if they are an adult
\emph{and} have a ticket'', ``if it is Friday \emph{or} a public holiday''.
Salam has three operators for combining conditions, and all three are plain
words:

\begin{center}
\begin{tabular}{cll}
\toprule
Operator & Meaning & The result is true when \\
\midrule
\code{and} & and & both sides are true \\
\code{or} & or & at least one of the two sides is true \\
\code{not} & negation & the side to its right is false \\
\bottomrule
\end{tabular}
\end{center}

\program{Combining conditions}
\begin{salamcode}
func main:
    age := 20
    has ticket := false
    println (age >= 18) and has ticket
    println (age >= 18) or has ticket
    println not has ticket
    println (age > 12) and (age < 65)
end
\end{salamcode}

\begin{salamout}
false
true
true
true
\end{salamout}

Look at the last line: to ask ``is the age between 12 and 65?'' we had to
join two separate comparisons with \code{and}. Writing \code{12 < age < 65},
the way we do in mathematics, means nothing in Salam.

\begin{note}
The signs \code{\&\&}, \code{||} and \code{!} used in many other languages
are errors in Salam; always write \code{and}, \code{or} and \code{not}.
\end{note}

\Needspace{16\baselineskip}
\begin{deep}[Truth tables]
If you would like to see every case of \code{and} and \code{or} in one
place, the table below will help. Do not memorise it; just compare it with
the everyday meaning of ``and'' and ``or'' and you will see that it is
exactly the same.
\begin{center}
\begin{tabular}{cccc}
\toprule
a & b & \code{a and b} & \code{a or b} \\
\midrule
true & true & true & true \\
true & false & false & true \\
false & true & false & true \\
false & false & false & false \\
\bottomrule
\end{tabular}
\end{center}
The only place where the programmer's ``or'' differs from the everyday one
is when both sides are true: in conversation, ``tea or coffee?'' means one
of the two, but \code{or} in Salam is also true when both are true.
\end{deep}

\section{A more complete program}

Suppose you want to know how many crates of drinks to buy for a party. Each
crate holds 6 bottles and each guest drinks 2:

\program{Shopping for a party}
\begin{salamcode}
func main:
    guests := 17
    bottles per guest := 2
    bottles per crate := 6

    bottles needed := guests * bottles per guest
    full crates := bottles needed / bottles per crate
    extra bottles := bottles needed % bottles per crate
    // if there are extra bottles, one more crate is needed
    crates needed := full crates + (extra bottles > 0 ? 1 : 0)

    println "Bottles needed:", bottles needed
    println "Crates needed:", crates needed
    println "Bottles left over:", (crates needed * bottles per crate) - bottles needed
end
\end{salamcode}

\begin{salamout}
Bottles needed: 34
Crates needed: 6
Bottles left over: 2
\end{salamout}

There is a new operator in this program: \code{condition ? a : b}. It says
``if the condition is true, the value a, otherwise the value b''. Here, if
any extra bottles are left, we buy one more crate. (You may wonder why we did
not call the third variable \code{bottles in crate}; the answer is in
chapter \ref{ch:variables}: \code{in} is one of the words of the language
and cannot be part of a name.) The next chapter, which is about making
decisions, has more to say about this operator.

\begin{summary}
\begin{itemize}
  \item The arithmetic operators are \code{+ - * / \%} and the power
        \code{**}. Integer division throws away the fraction, and
        \code{\%} gives the remainder.
  \item Order: powers first, then multiplication and division, then
        addition and subtraction. Use parentheses to change the order.
  \item Short forms: \code{+=}, \code{-=}, \code{*=}, \code{/=},
        \code{\%=}, \code{++}, \code{--}.
  \item Comparisons: \code{== != < > <= >=}, which give boolean results.
  \item Conditions are combined with \code{and}, \code{or} and \code{not}.
  \item \code{condition ? a : b} chooses one of two values according to the
        condition.
\end{itemize}
\end{summary}

\section*{Exercises}
\markright{Exercises}

\exercise Without running them, give the result of each expression, then
check: \code{5 + 2 * 3}, \code{(5 + 2) * 3}, \code{20 / 3 * 3},
\code{20 \% 3}, \code{2 ** 2 ** 3}.

\exercise Put a four-digit number into a variable and, using division and
remainder, work out and print the sum of its digits.

\exercise Put the price of an item and the amount the customer paid into two
variables, and split the change into the fewest possible 20, 10 and 5 dollar
notes, printing how many of each are needed.

\exercise Put someone's year of birth into a variable and, with a single
boolean expression, print whether they reach voting age (18) this year
(2026).

\exercise Make three boolean variables called \code{it is raining},
\code{i have an umbrella} and \code{i have a car}, and write an expression
that says whether you can go out ``without getting wet'': that is, either it
is not raining, or you have an umbrella or a car. Try the program with
several combinations of values.
